% !TEX program = pdflatex
% ==========================================================
% CLASE BEAMER PREMIUM
% Fundamentos de Matemática Básica - Agronomía
% Tema: Potencias
% Universidad Santo Tomás - Talca
% ==========================================================

\newif\ifhandout
%\handoutfalse
\handouttrue

\ifhandout
\documentclass[aspectratio=169,10pt,handout]{beamer}
\else
\documentclass[aspectratio=169,10pt]{beamer}
\fi

% ----------------------------------------------------------
% PAQUETES
% ----------------------------------------------------------
\usepackage[spanish,es-noshorthands]{babel}
\usepackage[T1]{fontenc}
\usepackage[utf8]{inputenc}
\usepackage{lmodern}
\usepackage{amsmath,amsfonts,amssymb}
\usepackage{ragged2e}
\usepackage{enumitem}
\usepackage{multicol}
\usepackage{tikz}
\usetikzlibrary{positioning,calc}
\usepackage{fontawesome5}
\usepackage{array}

% ----------------------------------------------------------
% COLORES
% ----------------------------------------------------------
\definecolor{USTBlue}{RGB}{0,70,127}
\definecolor{USTBlueDark}{RGB}{0,45,85}
\definecolor{USTTeal}{RGB}{0,153,117}
\definecolor{USTGray}{RGB}{120,120,120}
\definecolor{USTSoft}{RGB}{248,249,250}
\definecolor{USTText}{RGB}{35,40,45}
\definecolor{USTGold}{RGB}{196,154,70}
\definecolor{USTLightBlue}{RGB}{227,239,248}
\definecolor{USTLightTeal}{RGB}{226,245,240}
\definecolor{USTRedSoft}{RGB}{252,235,235}

% ----------------------------------------------------------
% TEMA
% ----------------------------------------------------------
\usetheme{Madrid}
\usecolortheme[named=USTBlue]{structure}
\setbeamertemplate{navigation symbols}{}

% ----------------------------------------------------------
% ESPACIADOS / LISTAS
% ----------------------------------------------------------
\setlist[itemize]{itemsep=0.35em}
\setlist[enumerate]{itemsep=0.35em}

% ----------------------------------------------------------
% COLORES GENERALES BEAMER
% ----------------------------------------------------------
\setbeamercolor{normal text}{fg=USTText,bg=white}
\setbeamercolor{frametitle}{fg=white,bg=USTBlueDark}
\setbeamercolor{title}{fg=USTBlueDark}
\setbeamercolor{subtitle}{fg=USTTeal}

\setbeamercolor{block title}{bg=USTLightBlue,fg=USTBlueDark}
\setbeamercolor{block body}{bg=USTSoft,fg=black}

\setbeamercolor{block title example}{bg=USTLightTeal,fg=USTBlueDark}
\setbeamercolor{block body example}{bg=USTSoft,fg=black}

\setbeamercolor{block title alerted}{bg=USTRedSoft,fg=black}
\setbeamercolor{block body alerted}{bg=red!2,fg=black}

% ----------------------------------------------------------
% FRAMETITLE PREMIUM
% ----------------------------------------------------------
\setbeamertemplate{frametitle}{
	\vspace*{-0.05cm}
	\begin{beamercolorbox}[wd=\paperwidth,ht=0.95cm,dp=0.18cm,leftskip=0.45cm,rightskip=0.25cm]{frametitle}
		{\large\bfseries \insertframetitle}
	\end{beamercolorbox}
}

% ----------------------------------------------------------
% PIE DE PÁGINA PREMIUM
% ----------------------------------------------------------
\setbeamercolor{author in head/foot}{bg=USTSoft,fg=USTGray}
\setbeamercolor{date in head/foot}{bg=USTSoft,fg=USTGray}
\setbeamercolor{title in head/foot}{bg=USTSoft,fg=USTGray}

\setbeamertemplate{footline}{%
	\leavevmode%
	\hbox{%
		\begin{beamercolorbox}[wd=.82\paperwidth,ht=2.8ex,dp=1ex,leftskip=0.4cm]{author in head/foot}%
			\scriptsize
			\textbf{Fundamentos de Matemática Básica}
			\quad|\quad
			Agronomía
			\quad|\quad
			Universidad Santo Tomás Talca
		\end{beamercolorbox}%
		\begin{beamercolorbox}[wd=.18\paperwidth,ht=2.8ex,dp=1ex,rightskip=0.35cm plus1fil]{date in head/foot}%
			\scriptsize \insertframenumber/\inserttotalframenumber
		\end{beamercolorbox}%
	}%
}

% ----------------------------------------------------------
% COMANDOS
% ----------------------------------------------------------
\newcommand{\destacar}[1]{\textcolor{USTBlueDark}{\textbf{#1}}}
\newcommand{\alerta}[1]{\textcolor{red!70!black}{\textbf{#1}}}
\newcommand{\pp}{\pause}

\newcommand{\iconoInicio}{\faPlayCircle}
\newcommand{\iconoDesarrollo}{\faBrain}
\newcommand{\iconoPractica}{\faPenSquare}
\newcommand{\iconoAgro}{\faSeedling}
\newcommand{\iconoCierre}{\faCheckCircle}
\newcommand{\iconoIdea}{\faLightbulb}
\newcommand{\iconoError}{\faExclamationTriangle}

\newenvironment{metaBox}
{\begin{block}{\faBullseye\ \ Meta de aprendizaje}}
	{\end{block}}

\newenvironment{ideaBox}
{\begin{exampleblock}{\iconoIdea\ \ Idea clave}}
	{\end{exampleblock}}

% ----------------------------------------------------------
% PORTADA PREMIUM
% ----------------------------------------------------------
\newcommand{\PortadaProfesionalUST}{%
	\begin{frame}[plain]
		\begin{tikzpicture}[remember picture,overlay]
			
			% Fondo
			\fill[white] (current page.south west) rectangle (current page.north east);
			
			% Franja superior
			\fill[USTBlueDark] (current page.north west) rectangle ([yshift=-1.45cm]current page.north east);
			\fill[USTTeal] ([yshift=-1.45cm]current page.north west) rectangle ([yshift=-1.62cm]current page.north east);
			
			% Banda lateral
			\fill[USTSoft] (current page.south west) rectangle ([xshift=2.3cm]current page.north west);
			
			% Motivos decorativos
			\fill[USTBlue!7] ([xshift=1.15cm,yshift=-2.0cm]current page.north west) circle (0.75cm);
			\fill[USTTeal!10] ([xshift=1.75cm,yshift=-4.1cm]current page.north west) circle (0.48cm);
			\fill[USTGold!15] ([xshift=-1.25cm,yshift=1.05cm]current page.south east) circle (1.10cm);
			
			% Caja principal
			\fill[black!7,rounded corners=8pt]
			([xshift=-5.1cm,yshift=-3.25cm]current page.center) rectangle
			([xshift=5.0cm,yshift=2.65cm]current page.center);
			
			\fill[white,rounded corners=8pt]
			([xshift=-5.25cm,yshift=-3.10cm]current page.center) rectangle
			([xshift=4.85cm,yshift=2.80cm]current page.center);
			
			% Barra lateral caja
			\fill[USTTeal] ([xshift=-5.25cm,yshift=-3.10cm]current page.center) rectangle
			([xshift=-4.95cm,yshift=2.80cm]current page.center);
			
			% Marca superior
			\node[anchor=north east,text=white]
			at ([xshift=-0.7cm,yshift=-0.52cm]current page.north east)
			{\small\bfseries Universidad Santo Tomás \quad Ciencias Básicas Talca};
			
			% Contenido central
			\node[align=center] at ([xshift=0.12cm,yshift=-0.08cm]current page.center){%
				\begin{minipage}{0.68\paperwidth}
					\centering
					
					{\Large\bfseries\color{USTText}\insertauthor\par}
					
					\vspace{0.7cm}
					
					{\fontsize{24}{28}\selectfont\bfseries\color{USTBlueDark}\inserttitle\par}
					
					\vspace{0.28cm}
					
					{\large\bfseries\color{USTTeal}\insertsubtitle\par}
					
					\vspace{0.55cm}
					
					{\normalsize\color{USTGray}\faSeedling\ Agronomía \quad \faUniversity\ Universidad Santo Tomás\par}
					
					\vspace{0.65cm}
					
					{\small\color{USTGray}\insertdate\par}
				\end{minipage}
			};
			
			% Línea inferior
			\draw[USTGray!35,line width=0.35pt]
			([xshift=0.9cm,yshift=0.72cm]current page.south west) --
			([xshift=-0.9cm,yshift=0.72cm]current page.south east);
			
			% Pie inferior
			\node[anchor=south,text=USTGray]
			at ([yshift=0.2cm]current page.south)
			{\scriptsize
				\textbf{Fundamentos de Matemática Básica}
				\;|\;
				Clase premium
				\;|\;
				Potencias};
			
		\end{tikzpicture}
	\end{frame}
}

% ----------------------------------------------------------
% DIAPOSITIVA DE SECCIÓN
% ----------------------------------------------------------
\newcommand{\SectionFrame}[3]{%
	\begin{frame}[plain]
		\begin{tikzpicture}[remember picture,overlay]
			\fill[USTBlueDark] (current page.south west) rectangle (current page.north east);
			\fill[USTTeal] (current page.south west) rectangle ([yshift=0.25cm]current page.north east);
			\fill[white,opacity=0.08] ([xshift=-1cm,yshift=1cm]current page.south east) circle (1.8cm);
			\fill[white,opacity=0.06] ([xshift=1.5cm,yshift=-1cm]current page.north west) circle (1.3cm);
			
			\node[anchor=center,text=white] at ([yshift=0.85cm]current page.center)
			{{\fontsize{34}{36}\selectfont #1}};
			
			\node[anchor=center,text=white] at (current page.center)
			{{\LARGE\bfseries #2}};
			
			\node[anchor=center,text=white!90] at ([yshift=-1.0cm]current page.center)
			{{\large #3}};
		\end{tikzpicture}
	\end{frame}
}

% ----------------------------------------------------------
% DATOS
% ----------------------------------------------------------
\title{Fundamentos de Matemática Básica}
\subtitle{Clase premium: Potencias}
\author{Prof. Luis González Alcaino}
\institute{Universidad Santo Tomás}
\date{Año académico 2026}

% ==========================================================
% DOCUMENTO
% ==========================================================
\begin{document}
	
	\PortadaProfesionalUST
	
	% ----------------------------------------------------------
	\begin{frame}{Panorama de la sesión}
		\begin{block}{\faListOl\ \ Organización de la clase (80 minutos)}
			\begin{enumerate}
				\item \destacar{Inicio:} activación de conocimientos previos.
				\item \destacar{Desarrollo:} definición y propiedades de las potencias.
				\item \destacar{Práctica guiada:} ejercicios resueltos paso a paso.
				\item \destacar{Práctica autónoma:} aplicación individual.
				\item \destacar{Cierre:} síntesis y reflexión.
			\end{enumerate}
		\end{block}
		
		\begin{metaBox}
			Aplicar el concepto de potencia y sus propiedades en ejercicios matemáticos y problemas contextualizados en Agronomía.
		\end{metaBox}
	\end{frame}
	
	% ----------------------------------------------------------
	\SectionFrame{\iconoInicio}{Inicio}{Activación y conexión con conocimientos previos}
	
	% ----------------------------------------------------------
	\begin{frame}{\iconoInicio\ \ Activación inicial}
		\begin{block}{Preguntas de entrada}
			\begin{itemize}
				\item ¿Cómo escribirías \(2\cdot2\cdot2\cdot2\) de forma abreviada?
				\item ¿Qué significa el exponente en una potencia?
				\item ¿Dónde aparecen multiplicaciones repetidas en contextos reales?
			\end{itemize}
		\end{block}
		
		\pp
		
		\begin{ideaBox}
			Las potencias permiten expresar de manera compacta una multiplicación repetida.
		\end{ideaBox}
	\end{frame}
	
	% ----------------------------------------------------------
	\begin{frame}{\iconoInicio\ \ ¿Por qué estudiar potencias en Agronomía?}
		\begin{columns}[T,totalwidth=\textwidth]
			\begin{column}{0.49\textwidth}
				\begin{block}{Aplicaciones}
					\begin{itemize}
						\item Densidad de siembra.
						\item Crecimiento de poblaciones bacterianas.
						\item Concentraciones muy pequeñas.
						\item Escalas técnicas y potencias de 10.
					\end{itemize}
				\end{block}
			\end{column}
			\begin{column}{0.47\textwidth}
				\begin{exampleblock}{Ejemplo técnico}
					\[
					0{,}00045\ \text{kg} = 4{,}5\times10^{-4}\ \text{kg}
					\]
					Esta forma facilita comparar, calcular y comunicar resultados.
				\end{exampleblock}
			\end{column}
		\end{columns}
	\end{frame}
	
	% ----------------------------------------------------------
	\SectionFrame{\iconoDesarrollo}{Desarrollo}{Concepto, propiedades y patrones}
	
	% ----------------------------------------------------------
	\begin{frame}{\iconoDesarrollo\ \ Concepto de potencia}
		\begin{block}{Definición}
			Una \destacar{potencia} representa una multiplicación repetida de un mismo número:
			\[
			a^n=\underbrace{a\cdot a\cdot a\cdots a}_{n\text{ veces}}
			\]
			donde:
			\begin{itemize}
				\item \(a\) es la \destacar{base},
				\item \(n\) es el \destacar{exponente}.
			\end{itemize}
		\end{block}
		
		\pp
		
		\begin{exampleblock}{Ejemplos}
			\[
			2^4=2\cdot2\cdot2\cdot2=16
			\qquad
			3^3=27
			\]
			\[
			5^2=25
			\qquad
			10^3=1000
			\]
		\end{exampleblock}
	\end{frame}
	
	% ----------------------------------------------------------
	\begin{frame}{\iconoDesarrollo\ \ Base y exponente}
		\begin{columns}[T,totalwidth=\textwidth]
			\begin{column}{0.48\textwidth}
				\begin{block}{Base}
					Es el número que se repite como factor.
					
					En \(5^3\), la base es \(5\).
				\end{block}
			\end{column}
			\begin{column}{0.48\textwidth}
				\begin{exampleblock}{Exponente}
					Indica cuántas veces se repite la base.
					
					En \(5^3\), el exponente es \(3\).
				\end{exampleblock}
			\end{column}
		\end{columns}
		
		\vspace{0.3cm}
		\centering
		\[
		5^3=5\cdot5\cdot5=125
		\]
	\end{frame}
	
	% ----------------------------------------------------------
	\begin{frame}{\iconoDesarrollo\ \ Propiedad 1: producto de igual base}
		\begin{block}{Regla}
			\[
			a^m\cdot a^n = a^{m+n}
			\]
		\end{block}
		
		\begin{exampleblock}{Ejemplo}
			\[
			2^3\cdot2^4 = 2^{3+4}=2^7=128
			\]
		\end{exampleblock}
		
		\begin{alertblock}{Interpretación}
			Si la base es la misma, \destacar{se suman los exponentes}.
		\end{alertblock}
	\end{frame}
	
	% ----------------------------------------------------------
	\begin{frame}{\iconoDesarrollo\ \ Propiedad 2: cociente de igual base}
		\begin{block}{Regla}
			\[
			\frac{a^m}{a^n}=a^{m-n}, \qquad a\neq 0
			\]
		\end{block}
		
		\begin{exampleblock}{Ejemplo}
			\[
			\frac{5^6}{5^2}=5^{6-2}=5^4=625
			\]
		\end{exampleblock}
		
		\begin{alertblock}{Interpretación}
			Si la base es la misma, \destacar{se restan los exponentes}.
		\end{alertblock}
	\end{frame}
	
	% ----------------------------------------------------------
	\begin{frame}{\iconoDesarrollo\ \ Propiedad 3: potencia de una potencia}
		\begin{block}{Regla}
			\[
			(a^m)^n = a^{m\cdot n}
			\]
		\end{block}
		
		\begin{exampleblock}{Ejemplo}
			\[
			(3^2)^4 = 3^{2\cdot4}=3^8=6561
			\]
		\end{exampleblock}
		
		\begin{alertblock}{Interpretación}
			En una potencia elevada a otra potencia, \destacar{se multiplican los exponentes}.
		\end{alertblock}
	\end{frame}
	
	% ----------------------------------------------------------
	\begin{frame}{\iconoDesarrollo\ \ Propiedad 4: potencia de un producto}
		\begin{block}{Regla}
			\[
			(ab)^n = a^n b^n
			\]
		\end{block}
		
		\begin{exampleblock}{Ejemplo}
			\[
			(2\cdot5)^3 = 2^3\cdot5^3 = 8\cdot125 = 1000
			\]
		\end{exampleblock}
		
		\begin{alertblock}{Interpretación}
			La potencia puede distribuirse a cada factor del producto.
		\end{alertblock}
	\end{frame}
	
	% ----------------------------------------------------------
	\begin{frame}{\iconoDesarrollo\ \ Propiedad 5: potencia de un cociente}
		\begin{block}{Regla}
			\[
			\left(\frac{a}{b}\right)^n = \frac{a^n}{b^n}, \qquad b\neq 0
			\]
		\end{block}
		
		\begin{exampleblock}{Ejemplo}
			\[
			\left(\frac{2}{3}\right)^3 = \frac{2^3}{3^3}=\frac{8}{27}
			\]
		\end{exampleblock}
	\end{frame}
	
	% ----------------------------------------------------------
	\begin{frame}{\iconoDesarrollo\ \ Propiedad 6: exponente cero}
		\begin{block}{Regla}
			\[
			a^0=1,\qquad a\neq0
			\]
		\end{block}
		
		\begin{exampleblock}{Ejemplo}
			\[
			7^0=1
			\]
		\end{exampleblock}
		
		\begin{alertblock}{Observación}
			Toda base distinta de cero elevada a cero vale 1.
		\end{alertblock}
	\end{frame}
	
	% ----------------------------------------------------------
	\begin{frame}{\iconoDesarrollo\ \ Propiedad 7: exponente negativo}
		\begin{block}{Regla}
			\[
			a^{-n}=\frac{1}{a^n},\qquad a\neq0
			\]
		\end{block}
		
		\begin{exampleblock}{Ejemplo}
			\[
			2^{-3}=\frac{1}{2^3}=\frac{1}{8}
			\]
		\end{exampleblock}
	\end{frame}
	
	% ----------------------------------------------------------
	\begin{frame}{\iconoDesarrollo\ \ Propiedad 8: casos especiales}
		\begin{block}{Reglas}
			\[
			a^1=a
			\qquad
			1^n=1
			\qquad
			0^n=0 \text{ para } n>0
			\]
		\end{block}
		
		\begin{exampleblock}{Ejemplos}
			\[
			9^1=9
			\qquad
			1^{12}=1
			\qquad
			0^5=0
			\]
		\end{exampleblock}
	\end{frame}
	
	% ----------------------------------------------------------
	\begin{frame}{\iconoError\ \ Errores frecuentes}
		\begin{alertblock}{Error 1}
			\[
			a^m+a^n \neq a^{m+n}
			\]
		\end{alertblock}
		
		\begin{exampleblock}{Ejemplo}
			\[
			2^2+2^3=4+8=12
			\qquad \text{pero} \qquad
			2^{2+3}=2^5=32
			\]
		\end{exampleblock}
		
		\pp
		
		\begin{alertblock}{Error 2}
			\[
			(a+b)^n \neq a^n+b^n
			\]
		\end{alertblock}
		
		\begin{exampleblock}{Ejemplo}
			\[
			(2+3)^2=25
			\qquad \text{pero} \qquad
			2^2+3^2=13
			\]
		\end{exampleblock}
	\end{frame}
	
	% ----------------------------------------------------------
	\begin{frame}{\iconoDesarrollo\ \ Patrones numéricos}
		\begin{columns}[T,totalwidth=\textwidth]
			\begin{column}{0.48\textwidth}
				\begin{block}{Potencias de 2}
					\[
					2^1=2,\quad 2^2=4,\quad 2^3=8
					\]
					\[
					2^4=16,\quad 2^5=32
					\]
				\end{block}
			\end{column}
			\begin{column}{0.48\textwidth}
				\begin{exampleblock}{Potencias de 10}
					\[
					10^1=10,\quad 10^2=100
					\]
					\[
					10^3=1000,\quad 10^4=10000
					\]
				\end{exampleblock}
			\end{column}
		\end{columns}
		
		\vspace{0.2cm}
		\begin{ideaBox}
			Cuando el exponente aumenta en 1, la potencia se multiplica por la base.
		\end{ideaBox}
	\end{frame}
	
	% ----------------------------------------------------------
	\SectionFrame{\iconoPractica}{Práctica guiada}{Resolución paso a paso}
	
	% ----------------------------------------------------------
	\begin{frame}{\iconoPractica\ \ Ejercicio guiado 1}
		\begin{block}{Problema}
			Simplificar:
			\[
			2^4\cdot2^3
			\]
		\end{block}
		
		\pp
		
		\begin{exampleblock}{Paso 1}
			La base es la misma, por lo tanto:
			\[
			2^4\cdot2^3 = 2^{4+3}
			\]
		\end{exampleblock}
		
		\pp
		
		\begin{exampleblock}{Paso 2}
			\[
			2^{4+3}=2^7=128
			\]
		\end{exampleblock}
	\end{frame}
	
	% ----------------------------------------------------------
	\begin{frame}{\iconoPractica\ \ Ejercicio guiado 2}
		\begin{block}{Problema}
			Resolver:
			\[
			(3^2)^3
			\]
		\end{block}
		
		\pp
		
		\begin{exampleblock}{Paso 1}
			Es una potencia de una potencia:
			\[
			(3^2)^3=3^{2\cdot3}
			\]
		\end{exampleblock}
		
		\pp
		
		\begin{exampleblock}{Paso 2}
			\[
			3^{2\cdot3}=3^6=729
			\]
		\end{exampleblock}
	\end{frame}
	
	% ----------------------------------------------------------
	\SectionFrame{\iconoAgro}{Aplicaciones agronómicas}{Matemática conectada con el contexto profesional}
	
	% ----------------------------------------------------------
	\begin{frame}{\iconoAgro\ \ Aplicación 1: semillas por superficie}
		\begin{block}{Situación}
			En una parcela experimental se utilizan:
			\[
			2^5
			\]
			semillas por metro cuadrado.
			
			Si la superficie cultivada es de:
			\[
			2^3
			\]
			metros cuadrados, determine el total de semillas.
		\end{block}
	\end{frame}
	
	% ----------------------------------------------------------
	\begin{frame}{\iconoAgro\ \ Resolución de la aplicación 1}
		\begin{exampleblock}{Paso 1}
			Multiplicamos:
			\[
			2^5\cdot2^3
			\]
		\end{exampleblock}
		
		\pp
		
		\begin{exampleblock}{Paso 2}
			Como la base es la misma:
			\[
			2^5\cdot2^3=2^{5+3}=2^8
			\]
		\end{exampleblock}
		
		\pp
		
		\begin{exampleblock}{Paso 3}
			\[
			2^8=256
			\]
			Se utilizan \textbf{256 semillas}.
		\end{exampleblock}
	\end{frame}
	
	% ----------------------------------------------------------
	\begin{frame}{\iconoAgro\ \ Aplicación 2: crecimiento bacteriano}
		\begin{block}{Situación}
			Una colonia bacteriana del suelo se duplica cada hora.
			
			Si inicialmente hay:
			\[
			2^3
			\]
			bacterias, después de 4 horas habrá:
			\[
			2^3\cdot2^4
			\]
		\end{block}
		
		\pp
		
		\begin{exampleblock}{Resolución}
			\[
			2^3\cdot2^4 = 2^{7}=128
			\]
			Después de 4 horas hay \textbf{128 bacterias}.
		\end{exampleblock}
	\end{frame}
	
	% ----------------------------------------------------------
	\begin{frame}{\iconoAgro\ \ Aplicación 3: fertilización}
		\begin{block}{Situación}
			Un fertilizante se aplica en dosis de:
			\[
			5\times10^2 \text{ gramos por hectárea}
			\]
			Si el ensayo abarca:
			\[
			10^3 \text{ hectáreas}
			\]
			¿cuántos gramos se necesitan en total?
		\end{block}
		
		\pp
		
		\begin{exampleblock}{Resolución}
			\[
			(5\times10^2)(10^3)=5\times10^5
			\]
			\[
			5\times10^5=500000
			\]
			Se requieren \textbf{500000 gramos}.
		\end{exampleblock}
	\end{frame}
	
	% ----------------------------------------------------------
	\begin{frame}{\iconoAgro\ \ Aplicación 4: micronutrientes}
		\begin{block}{Situación}
			Una solución nutritiva contiene:
			\[
			3\times10^{-2}\ \text{g}
			\]
			de un micronutriente por litro.
			
			Si se preparan:
			\[
			10^3
			\]
			litros, ¿cuántos gramos se requieren?
		\end{block}
		
		\pp
		
		\begin{exampleblock}{Resolución}
			\[
			(3\times10^{-2})(10^3)=3\times10^{1}=30
			\]
			Se necesitan \textbf{30 gramos}.
		\end{exampleblock}
	\end{frame}
	
	% ----------------------------------------------------------
	\SectionFrame{\iconoPractica}{Práctica autónoma}{Trabajo individual y consolidación}
	
	% ----------------------------------------------------------
	\begin{frame}{\iconoPractica\ \ Ejercicios propuestos}
		Resuelva individualmente:
		\begin{enumerate}
			\item \(2^5\cdot2^2\)
			\item \(\dfrac{3^7}{3^4}\)
			\item \((4^2)^3\)
			\item \((2\cdot3)^2\)
			\item \(\left(\dfrac{3}{5}\right)^2\)
			\item \(8^0\)
			\item \(2^{-4}\)
			\item En un cultivo hay \(3^4\) plantas por sector y \(3^2\) sectores. ¿Cuántas plantas hay en total?
		\end{enumerate}
	\end{frame}
	
	% ----------------------------------------------------------
	\begin{frame}{\iconoPractica\ \ Respuestas}
		\begin{multicols}{2}
			\begin{enumerate}
				\item \(2^5\cdot2^2=2^7=128\)
				\item \(\dfrac{3^7}{3^4}=3^3=27\)
				\item \((4^2)^3=4^6=4096\)
				\item \((2\cdot3)^2=2^2\cdot3^2=36\)
				\item \(\left(\dfrac{3}{5}\right)^2=\dfrac{9}{25}\)
				\item \(8^0=1\)
				\item \(2^{-4}=\dfrac{1}{16}\)
				\item \(3^4\cdot3^2=3^6=729\)
			\end{enumerate}
		\end{multicols}
	\end{frame}
	
	% ----------------------------------------------------------
	\SectionFrame{\iconoCierre}{Cierre}{Síntesis, reflexión y salida}
	
	% ----------------------------------------------------------
	\begin{frame}{\iconoCierre\ \ Síntesis de la clase}
		\begin{block}{Ideas clave}
			\begin{itemize}
				\item Una potencia representa una multiplicación repetida.
				\item La base es el número que se repite.
				\item El exponente indica cuántas veces se repite.
				\item Las propiedades permiten simplificar cálculos.
				\item Las potencias modelan fenómenos relevantes en Agronomía.
			\end{itemize}
		\end{block}
		
		\begin{ideaBox}
			Comprender potencias ayuda no solo a calcular mejor, sino también a interpretar situaciones reales del ámbito agrícola.
		\end{ideaBox}
	\end{frame}
	
	% ----------------------------------------------------------
	\begin{frame}{\iconoCierre\ \ Pregunta de salida}
		\begin{exampleblock}{Reflexión final}
			¿Por qué las potencias permiten representar de manera eficiente fenómenos de crecimiento y escalas en contextos agrícolas?
		\end{exampleblock}
		
		\vspace{0.4cm}
		
		\begin{center}
			{\Large\color{USTTeal}\faSeedling}
			
			\vspace{0.25cm}
			
			{\normalsize\textcolor{USTGray}{La matemática básica es una herramienta para comprender y tomar decisiones técnicas con mayor precisión.}}
		\end{center}
	\end{frame}
	
\end{document}